\section{한계}

\begin{enumerate}
  \item \textbf{수렴 차수를 보장하지 않음.} 확률적 점프, 벡터가 아닌 점의 외삽, 재귀적 이력,
  매끄럽지 않은 사영 때문에 AB2의 차수 논증을 직접 적용할 수 없다.
  \item \textbf{경계 기하.} 정확히 0인 성분은 Fisher--Rao 계량이 특이해지는 경계에 예측을 놓는다.
  \item \textbf{위치별 근사.} 각 토큰 위치를 서로 분리된 범주형 단체로 취급한다. 전체 결합 시퀀스
  분포의 기하는 표현하지 않는다.
\end{enumerate}

\section{결론}

본 논문은 Fisher--Rao 제곱근 구면 위에서 범주형 예측을 재귀적으로 외삽하는 MDLM 및 SEDD의
추론 시점 수정법인 측지 모멘텀 샘플링을 제시했다. 양의 직교영역 사영 이전의 갱신은 정확한
지수사상 해석을 가지며, 국소 계수는 가변 스텝 AB2 계수와 일치한다. 그러나 이 알고리즘은
확률흐름 ODE의 평행이동된 벡터장이 아니라 확률적 이산 경로에서 얻은 예측 점에 작용한다. 따라서
수학적으로 정확한 명칭은 증명된 2차 리만 솔버가 아니라 \textbf{재귀적 Fisher--Rao 구면 외삽}이다.

결과는 특히 MDLM의 적은 스텝에서 유망한 퍼플렉서티 경향을 보이지만, 분포 지표와 SEDD의
고스텝 결과는 혼재한다. 형식적인 수렴성을 주장하려면 연속 벡터장에 기반한 유도가 필요하다.
